\documentclass[tikz,border=6pt]{standalone}
\usepackage{ctex}
\usepackage{amsmath}
\usepackage{amssymb}
\usepackage{pgfplots}
\pgfplotsset{compat=1.18}
\usetikzlibrary{calc,angles,quotes,arrows.meta,positioning,decorations.markings}
\begin{document}
\begin{tikzpicture}

% ===== 左图：线性纵轴上 y=10^x 爆炸 =====
\begin{axis}[
    name=left,
    width=8cm, height=7.2cm,
    xlabel={$x$}, ylabel={$y$（线性刻度）},
    xmin=0, xmax=6, ymin=0, ymax=1100000,
    axis lines=left,
    xtick={0,1,2,3,4,5,6},
    ytick={0,200000,400000,600000,800000,1000000},
    yticklabels={$0$,$2{\times}10^5$,$4{\times}10^5$,$6{\times}10^5$,$8{\times}10^5$,$10^6$},
    scaled y ticks=false,
    every axis y label/.style={at={(ticklabel cs:1.02)},rotate=0,anchor=south},
    title={\bfseries 线性纵轴：$y=10^x$ 急剧爆炸},
    title style={yshift=2pt},
    tick label style={font=\footnotesize},
    label style={font=\small},
    enlarge x limits=false,
]
\addplot[blue,very thick,domain=0:6,samples=200] {10^x};
% 标注几个点
\addplot[only marks,mark=*,blue,mark size=1.6pt] coordinates {(1,10)(2,100)(3,1000)(4,10000)(5,100000)(6,1000000)};
\node[blue,font=\scriptsize,anchor=west] at (axis cs:6,1000000) {$(6,10^6)$};
\node[fill=white,fill opacity=0.85,text opacity=1,draw=blue!40,rounded corners,
      font=\scriptsize,align=left,inner sep=3pt] at (axis cs:2.05,720000)
  {低处的 $10,10^2,10^3$\\几乎贴在 $x$ 轴上，\\完全看不出差别};
\end{axis}

% ===== 右图：对数纵轴上 y=10^x 拉直成直线 =====
\begin{axis}[
    name=right,
    at={($(left.east)+(2.3cm,0)$)}, anchor=west,
    width=8cm, height=7.2cm,
    xlabel={$x$}, ylabel={$\lg y=\log_{10}y$（对数刻度）},
    xmin=0, xmax=6, ymin=0, ymax=6.4,
    axis lines=left,
    xtick={0,1,2,3,4,5,6},
    ytick={0,1,2,3,4,5,6},
    yticklabels={$10^0$,$10^1$,$10^2$,$10^3$,$10^4$,$10^5$,$10^6$},
    title={\bfseries 对数纵轴：取对数后\,拉直成直线},
    title style={yshift=2pt},
    tick label style={font=\footnotesize},
    label style={font=\small},
    enlarge x limits=false,
]
% 取对数后 lg(10^x)=x，是过原点斜率1的直线
\addplot[red,very thick,domain=0:6,samples=2] {x};
\addplot[only marks,mark=*,red,mark size=1.6pt] coordinates {(1,1)(2,2)(3,3)(4,4)(5,5)(6,6)};
\node[red,font=\scriptsize,anchor=north west] at (axis cs:4.4,4.05) {$\lg(10^x)=x$};
\node[fill=white,fill opacity=0.85,text opacity=1,draw=red!40,rounded corners,
      font=\scriptsize,align=left,inner sep=3pt] at (axis cs:1.75,5.15)
  {每个数量级\\等距排列：\\差距一目了然};
\end{axis}

\end{tikzpicture}
\end{document}
